% =====================================================================
% discussion.tex — upgraded empirical paper.
% THE model is affine_ew, reported throughout as "the model." No
% architecture comparison. Numbers taken verbatim from ACT_I_VDA_RESULTS
% and ACT_II_RESULTS (affine_ew values only). Arc: (1) directly measures
% criterion, sensitivity, and spatial gating; (2) the attention–
% working-memory link (read value, deploy, gate sensitivity, rehearse);
% (3) LIP-like priority map as an honest biological analogy; (4) forward
% pointers to the open sections; (5) honest limitations.
% =====================================================================

\section{Discussion}
\label{sec:discussion}

An affine-feedback recurrent vision-transformer architecture with a spatial
working-memory state, trained with reinforcement plus temporal JEPA self-supervision
(loss weight $0.5$) to detect a change among cued
Gabor stimuli, reproduces several functional signatures of primate visual attention
across a core checkpoint and separately trained task variants and --- through the attention-clamp battery --- takes several
of its long-standing open questions a step past what behaviour alone can
resolve. The model orients to a cue, relaxes near baseline during the blank and later
re-acquires cue-directed attention, gates its own detection sensitivity, shifts
its decision criterion, and holds a linearly decodable, position-specific memory signal
after the change. No cue-location decoder was fitted across the delay. These are
observed signatures of checkpoints optimized jointly by RL and a JEPA auxiliary loss
(weight $0.5$),
not hand-specified attention targets. We organise what follows around three claims the
results support --- that the model directly operationalizes criterion and sensitivity
effects plus spatial gating, that it exposes interactions between attention and working
memory, and that the joint architecture is usefully read against a
counterpart in the parietal priority map --- and then delimit, plainly,
where the account is incomplete.

\subsection*{Directly measured criterion, sensitivity, and spatial gating}

For twenty-five years the field has carried three competing explanations
of how attention improves perception: attention enhances the signal
(gain), it reduces external noise and selects the target more
efficiently, or it acts at the decision stage by shifting the observer's
criterion and reducing spatial uncertainty
\cite{carrasco2011_visual_attention_25y}. The three ``are not mutually exclusive;
rather, all three are likely to contribute'' \cite{carrasco2011_visual_attention_25y}.
The present battery does not estimate how much each contributes. Instead, it directly
measures criterion and sensitivity under the same attention clamp and tests whether an
injection effect is spatially gated. It contains no dedicated external-noise manipulation,
so efficient external-noise exclusion remains compatible with the results rather than
quantitatively isolated.

The spatial-gating result is direct:
attention added to the monitored patch raises detection, and injecting
priority onto a patch on no-change trials conjures a false detection ---
the false-alarm rate rises roughly fourfold, from about $0.03$ to about
$0.12$, as injection strength grows (Section~\ref{sec:detection}). This
is an intervention effect in the decision pathway. It is spatially gated:
injecting priority at the monitored location evokes the false
percept, whereas injecting at an unmonitored patch does essentially
nothing (a flat $\sim$$0.05$--$0.10$); the model reads out the location
it is already selecting, not priority in the abstract
(Section~\ref{sec:detection}). A cued-versus-uncued benefit is present in
VDA4, whereas VDA1 has no uncued active item and therefore does not define
the same spatial contrast. Separately, the low-to-high-validity
\emph{cued-threshold} changes are $0.148^\circ$, $0.193^\circ$,
$0.979^\circ$, and $2.207^\circ$ across VDA1/2/4/9. The $2.207^\circ$
endpoint is not a spatial cued-versus-uncued benefit. These associations are compatible with a noise-limited selector
\cite{carrasco2011_visual_attention_25y, MullerFindlay1987}, but VDA9 also
changes grid, token count, readout dimension, validity semantics, and
training support, and they are not an external-noise test. \emph{Decision criterion}: in the canonical VDA4
checkpoint, sweeping the nominal clamp from suppression to enhancement
moves the sampled signal-detection criteria in the liberal direction,
$c: 1.42 \to -0.60$ (Section~\ref{sec:detection}) --- the textbook
attentional criterion shift \cite{MullerFindlay1987,posner1980_orienting},
now produced by an analysis intervention rather than a payoff instruction.
This monotonic description is local to VDA4; VDA9 reverses at the
high-clamp end.

The criterion and sensitivity endpoints are both operationalized on the
\emph{same} clamp. The criterion moves smoothly across the VDA4 sweep;
sensitivity is load-bearing on attention ---
suppressing attention collapses $d'$ (to $0.63$/$1.45$ in the
suppression regime) while natural and enhanced attention restore it (to
$\sim$$2.0$--$2.7$) (Section~\ref{sec:detection}). These are two directly measured
responses to the intervention, not a quantitative three-way contribution decomposition.
They parallel the distinction the
normalization account of gain predicts at the neural level
\cite{reynolds_heeger2009_normalization, mcadams_maunsell1999_v4_tuning,
treue_martinez_trujillo1999_feature_attention} and which psychophysics
can only reach indirectly. The same contrast appears in a preliminary cued analogue of the
Luo--Maunsell paradigm: in separately trained reward-regime checkpoints, concentrating value at one
location moves sensitivity ($\Delta d' = 1.09$ at $\Delta = 18^{\circ}$,
with a low false-alarm rate of $0.02$) while a global reward that spares
no location moves the criterion instead (false-alarm rate $0.18$,
sensitivity change of only $0.59$) (Section~\ref{sec:luomaunsell})
\cite{luo_maunsell2018_criterion_sensitivity}. That the two levers separate under reward
structure, and that related quantities separate under a direct attention clamp in
the base task, is encouraging. It is not yet a faithful cue-free, shared-base
reproduction, and the between-checkpoint comparison should not be read as a
within-model double dissociation.

One further, novel read-out sharpens the criterion story into a failure
mode. When attention is misdirected onto an empty monitored patch while
a real change occurs elsewhere, the model misses the change --- the
declare rate falls from $0.68$ to $0.24$ --- and does so while becoming
\emph{more} confident, its policy entropy dropping from $0.053$ to
$0.035$ (Section~\ref{sec:detection}). Misallocated attention does not
merely lower performance; it manufactures confident perceptual errors,
because the detector reports ``no change'' precisely because it was
pointed at the wrong place. This is a mechanistic account, in one system,
of how an attentional misallocation --- the sort a misleading cue induces
--- can produce not just misses but assured misses.

\subsection*{One architecture links attention to working memory}

The second claim is that attention and working memory interact within one recurrent
architecture. The current observations establish coupling, not identity: the memory
state shapes attention, attention gates sensitivity, and linear probes recover selected
task variables from memory at sampled frames.

The model \emph{reads value}: cue colour, which signals the worth of a
correct detection, is latched into the memory state at cue onset and
held perfectly --- decoding balanced accuracy $1.00$ from the first
frame through to the change (Section~\ref{sec:detection}). It
\emph{deploys spatially}: attention orients to the cued patch at cue
onset, well above the uniform baseline, so that the memory of \emph{where}
to look is expressed as a priority map over the visual field
(Section~\ref{sec:detection}). It \emph{gates its own sensitivity}: as
shown above, detection $d'$ rises and falls with the attention placed on
the monitored location, so the deployed attention is not a spectator but
the quantity that sets how well the change can be seen. Its blank-period state is
compatible with a rehearsal hypothesis, but does not establish one. In the canonical
affine VDA4 trace, cued attention falls near the uniform baseline during the blank
($0.22$ versus $0.25$), and no delay-probe or causal delay-clamp result is reported
(Section~\ref{sec:rehearsal}). Whether covert spatial attention is the load-bearing
maintenance mechanism therefore remains a prospective test
\cite{awh_vogel_oh2006,baddeley1992working}. The later cue-directed re-acquisition is
compatible with an intervening spatial trace, but no cue-location decoder was fitted
across the blank.

What distinguishes reading value from perceiving predictiveness is
instructive. Cue \emph{validity} --- how reliably the cue predicts the
change location --- is linearly decodable at the cue but falls toward
chance within the next sampled steps (Section~\ref{sec:detection}). That
weak maintenance accompanies the nearly flat validity effect in the base
task, but decoding alone does not prove that loss of linear recoverability
is its unique cause. The separately trained Validity4 task shows that the
architecture can use validity when it is the only cue on offer. There the
checkpoint uses it gradedly: the cueing benefit grows monotonically with displayed
validity ($0.00 \to 0.07 \to 0.08 \to 0.12$ from $25\%$ to $100\%$),
uncued detection falls as the cue sharpens ($0.94 \to 0.86$), and
attention on the cued patch rises with validity ($\alpha_c: 0.27 \to
0.48$ against a uniform $0.25$) (Section~\ref{sec:value}). A Posner-style
validity effect \cite{posner1980_orienting} is thus within the
architecture's repertoire. Under the legacy Validity4 sampler, displayed $p=0.25$
realizes $0.4375$ cued-change probability; archived VDA4 uses the same four-location
rule. Target semantics therefore do not confound this pair, although the separately
trained checkpoints and the presence versus absence of a value cue do. Their
contrast is compatible with cue competition but does not
establish that value monopolizes a fixed resource or causes validity to
be discarded. In these checkpoints, value remains linearly recoverable
from cue onset, whereas validity maintenance depends on task and
checkpoint.

\subsection*{A biological anchor: the priority map that binds position memory to detection threshold}

These behaviours motivate a functional analogy to a parietal priority
map. After the change, the model combines a position-specific event-location signal with an
attentional modulation of detection threshold (the clamp moves $d'$ and
$c$). The current analyses do not directly quantify a trial-level coupling
between representational strength and threshold. They nevertheless place
both measurable functions within one recurrent architecture, paralleling
accounts of area LIP and allied parietal circuits in which spatial representations
track task-relevant locations and bias visual processing there \cite{bisley_goldberg2010}.

We make this correspondence as an \emph{honest analogy, not an
isomorphism}. The model's memory is a recurrent state shaped by joint RL and JEPA optimization
(auxiliary weight $0.5$);
LIP is a population of parietal neurons embedded in a
cortical--subcortical circuit, and no element of the model corresponds to
a specific cell class, layer, or projection. What transfers is the
proposed computation --- event-location coding and attentional threshold
modulation within one architecture --- not anatomy or an established
single-signal identity. A direct relationship between code strength and
threshold remains to be tested.

\subsection*{Completed battery and open tests}

The affine $d_{\mathrm{mem}}=128$ VDA1/2/4/9 battery is reported in full
in Section~\ref{sec:vda-battery}; it is not forthcoming. Each point uses
one canonical checkpoint with one 20{,}000-row logged phase. The battery
shows an ordered increase in low-to-high-validity cued-threshold change,
but sustained attention, criterion range, and validity maintenance do not
form a single monotone law. VDA16 is an interrupted run ending near chance
at checkpoint 599, not a completed endpoint or an exhausted-budget
failure.

Two interventions remain open. Preserved $d_{\mathrm{mem}}=128$ and $256$
checkpoints exist, but the matched held-out psychometric, decoding, and
clamp battery has not been completed; no width or capacity law is reported
(Section~\ref{sec:capacity}). The delay-probe environment is implemented,
but no trained probe checkpoint, memorized-location advantage, or causal
delay-clamp result is reported. Attention-based rehearsal therefore
remains a hypothesis (Section~\ref{sec:rehearsal}).

\subsection*{Scope and limitations}

Several boundaries delimit these statements. The strongest is a training
outcome, not an interpretive caveat: not every attentional paradigm we
attempted was learnable by the model as posed. An attend-here/ignore-here
filtering task did not train --- the policy collapsed into always
withholding, never declaring a change, despite an environment whose
reward structure is achievable in principle. This failed task is not reported as a
quantitative Results section; its run-level trace is identified in the Appendix. The collapse is a learning failure of a
particular curriculum, not evidence that the mechanism is absent, but it
means we cannot report the corresponding selective-attention result, and
we say so rather than presenting a partial run as a finding. Relatedly,
the Validity4 checkpoint uses cue validity when validity is the operative cue,
whereas VDA4 weakly maintains validity when a competing value cue is present
(Sections~\ref{sec:detection} and~\ref{sec:value}). Because these are separately
trained checkpoints, the contrast is a task association rather than a within-weight
test of cue competition and should not be read as a general Posner effect.

The set-size characterization is complete for the four included checkpoints but is
single-checkpoint and confounded by task/model differences. The matched width analysis
and delay-probe rehearsal result are not in hand, and nothing in the discussion should
be read as asserting their outcomes. The Baruni relative-value test is a principled
null rather than a positive result: the model's value sensitivity is absolute, not
relative to context (Section~\ref{sec:value}), which is consistent with the regime the
model occupies but is a null nonetheless \cite{baruni2015}. Finally, the model family
is trained on a narrow family of change-detection tasks. The behaviours it reproduces
are broad, but the reproductions are functional signatures --- psychometric,
signal-detection, and decoding profiles --- and the biological correspondences are
analogies at the level of computation. The contribution is that an instrumented
jointly RL--JEPA-trained architecture (auxiliary weight $0.5$), dissected with causal clamps and linear probes, separates
several co-present attentional signatures while exposing a working-memory state whose
relation to rehearsal remains a prospective causal test.

% FIGURE: discussion synthesis schematic — optional. One panel mapping
% each measured effect (gain: microstim false-alarm curve; selection:
% checkpoint-specific cued-vs-uncued thresholds; criterion & sensitivity: c(alpha)
% and d'(alpha) clamp curves; value/validity/change decoding vs timestep;
% rehearsal: hypothesis only -- affine VDA4 blank attention is near baseline,
% and no trained delay-probe checkpoint, enhancement, or causal delay-clamp
% result exists) onto the three
% Carrasco accounts and the priority-map triple. Source npz: vda_sweep/
% figs/{microstim_curves, sdt_curves, decode_curves, alpha1_lines}
% (affine_ew series only).
\label{fig:discussion-synthesis}
